V predloženej bakalárskej práci sa nám úspešne podarilo navrhnúť a implementovať funkčný model autonómneho vozidla s detekciou prekážok. Preskúmali sme pritom možnosti a obmedzenia platformy mikrokontrolérov Arduino a jej prepojenia so softvérom na počítačové videnie OpenCV. Pri samotnej implementácii riešenia sme skombinovali softvérové aj hardvérové prostriedky na detekciu prekážok.\par
V našej práci sme využívali rodinu mikrokontrolérov Arduino. Daný výrobca poskytuje nákladovo efektívne riešenia, avšak s nižším výpočtovým výkonom, preto sme preskúmali možnosť distribuovaného systému s architektúrou klient-server. Vďaka podpore bezdrôtovej komunikácie na doske Arduino Nano ESP32, ako aj na kamerovom module ESP32-CAM, sa nám podarilo vytvoriť jednoduchú sieť. Presunutím náročných výpočtov na server sme síce obetovali časť reakčnej rýchlosti, no získali sme vyššiu robustnosť systému. Na druhej strane je tento prístup limitovaný kvalitou pripojenia, a preto nie je vhodný pre bezpečnostno-kritické aplikácie.\par
Na autonómne riadenie sme použili algoritmus detekcie vozovky využívajúci funkcie z knižnice OpenCV. Tento prístup sa vyznačoval vysokou presnosťou na rovných úsekoch, avšak v zákrutách mohlo dôjsť k strate sledovanej dráhy. Tento nedostatok by sa dal optimalizovať pomocou kontextových pomôcok, ktoré by vozidlu signalizovali, aby spomalilo.\par
Na detekciu prekážok sme použili konvolučnú neurónovú sieť YOLOv8. Model S (small) sa ukázal ako ideálny pre aplikácie v reálnom čase s väčším výpočtovým výkonom, ktorým sme disponovali na laptope. Táto sieť bola schopná detekovať rôzne triedy objektov, čo bolo pre našu implementáciu dôležité, avšak pri malých a čiastočne prekrytých prekážkach narážala na problémy. Tento fakt vyplýva z obmedzeného rozlíšenia a úzkeho zorného poľa kamery. Integráciou viacerých kamier a dodatočným trénovaním neurónovej siete by však bolo možné tento problém úspešne eliminovať.\par
Sekundárnym systémom detekcie v našej implementácii bol ultrazvukový senzor. Jeho uplatnenie sa ukázalo ako veľmi presné a nákladovo efektívne v rámci jeho určeného využitia. Pri detekcii na krátke vzdialenosti vykazoval vysokú presnosť, avšak jeho efektívny detekčný kužeľ bol malý. V reálnej implementácii by preto bolo ideálne rozmiestniť viacero takýchto senzorov po celom obvode vozidla.\par
Naša implementácia sa zameriavala na otestovanie distribuovanej architektúry autonómneho riadenia. Hoci je znižovanie nákladov v oblasti autonómnych systémov výrazným trendom, nemožno opomenúť, že bezpečnosť je vždy prvoradá. Z tohto dôvodu sa domnievame, že takéto systémy nemajú v automobilovom priemysle priamu budúcnosť. Uplatnenie by však mohli nájsť (a čiastočne ho už nachádzajú) v oblasti autonómneho doručovania tovaru, kde sú priemerná rýchlosť vozidiel a miera nebezpečenstva pre okolie pomerne nízke. Presunutím výpočtovej záťaže z edge na cloud computing by bolo možné výrazne znížiť výrobné náklady takejto doručovacej jednotky.\par
Na záver možno konštatovať, že predložená práca úspešne demonštruje uplatnenie platformy Arduino v autonómnej mobilite. Text detailne poukazuje na možnosti a nedostatky hardvéru i softvéru pri detekcii prekážok a vozovky. Práca zároveň prináša cenné poznatky o limitoch a praktickom využití distribuovanej architektúry, čím otvára nové otázky o budúcnosti a smerovaní autonómnej dopravy.